% from ega2-1-fr.tex
\oldpage[II]{121}

ce qui s'écrit aussi
\[
\label{II.6.4.1.4-fr}
(\det u)\cdot I_E=uQ(u)
\tag{6.4.1.4}
\]
($I_E$ automorphisme identique de $E$), avec
\[
\label{II.6.4.1.5-fr}
Q(u)=(-1)^{n+1}(u^{n-1}-\sigma_1(u)u^{n-2}+\ldots+(-1)^{n-1}\sigma_{n-1}(u)).
\tag{6.4.1.5}
\]

Soit $\varphi:A\to B$ un homomorphisme d'anneaux, faisant de $B$ une
$A$-algèbre~; considérons le $B$-module
$E_{(B)}=E\otimes_AB$ qui est libre de rang $n$, et l'extension $u\otimes1$
de $u$ à un endomorphisme de $E_{(B)}$~; il est immédiat que l'on a
$\sigma_i(u\otimes1)=\varphi(\sigma_i(u))$ pour tout indice $i$.
\end{env}

\begin{env}[6.4.2]
\label{II.6.4.2-fr}
Supposons maintenant que $A$ soit un anneau \emph{intègre}, de corps des
fractions $K$, et $E$ un $A$-module de \emph{type fini} (mais non
nécessairement libre). Soit $n$ le \emph{rang} de $E$, c'est-à-dire le rang du
$K$-module libre $E\otimes_AK$~; à tout endomorphisme $u$ de $E$ correspond
canoniquement l'endomorphisme $u\otimes1$ de $E\otimes_AK$~; par abus de
langage, on appelle encore \emph{polynôme caractéristique} de $u$ et on note
$P(u,T)$ le polynôme $P(u\otimes1,T)$, dont les coefficients
$\sigma_i(u\otimes1)$ (qui appartiennent à $K$) seront encore notés
$\sigma_i(u)$ et appelés les \emph{fonctions symétriques élémentaires de
$u$}~; en particulier $\det u=\det(u\otimes1)$ par définition. Avec ces
notations, les formules (\hyperref[II.6.4.1.3-fr]{6.4.1.3}) à
(\hyperref[II.6.4.1.5-fr]{6.4.1.5}) ont un sens et sont encore valables,
lorsque l'on interprète $u^j$ comme l'homomorphisme $E\to E\otimes_AK$
composé de l'endomorphisme $u^j\otimes1=(u\otimes1)^j$ de $E\otimes_AK$ et
de l'homomorphisme canonique $x\mapsto x\otimes1$.

Si $F$ est le sous-module de torsion de $E$, et si $E_0=E/F$, on a
$u(F)\subset F$, donc, par passage aux quotients, $u$ donne un endomorphisme
$u_0$ de $E_0$~; en outre $E\otimes_AK$ s'identifie à $E_0\otimes_AK$ et
$u\otimes1$ à $u_0\otimes1$, donc $\sigma_i(u)=\sigma_i(u_0)$ pour
$1\leq i\leq n$.

Lorsque $E$ est \emph{sans torsion}, $E$ s'identifie canoniquement à un
sous-$A$-module de $E\otimes_AK$, et la relation $u\otimes1=0$ est
équivalente à $u=0$. Lorsque $E$ est un $A$-module libre, les deux définitions
des $\sigma_i(u)$ données dans (\hyperref[II.6.4.1-fr]{6.4.1}) et
(\hyperref[II.6.4.2-fr]{6.4.2}) coïncident d'après ce qui précède, ce qui
justifie les notations adoptées.

On notera aussi que lorsque $E$ est un \emph{module de torsion}, autrement dit
$E_0=\{0\}$, l'algèbre extérieure de $E_0$ est réduite à $K$ et le déterminant
de l'unique endomorphisme $u_0$ de $E_0$ est \emph{égal à $1$}.
\end{env}

\begin{proposition}[6.4.3]
\label{II.6.4.3-fr}
Soient $A$ un anneau intègre, $E$ un $A$-module de type fini, $u$ un
endomorphisme de $E$~; alors les fonctions symétriques élémentaires
$\sigma_i(u)$ de $u$ (et en particulier $\det u$) sont des éléments de $K$
entiers sur $A$.
\end{proposition}

Soit $A'$ la clôture intégrale de $A$~; comme $A'[T]$ est un anneau
intégralement clos ([24], p.~99), c'est la clôture intégrale de
$A[T]$ dans son corps des fractions $K(T)$. Remplaçant $u$ par
$T\cdot I-u\otimes1$ et $A$ par $A[T]$, on voit qu'on est ramené à prouver que
$\det u$ est entier sur $A$. Si $n$ est le rang de $E$, on a
$\det u=\det(\bigwedge^n u)$ et
$(\bigwedge^n u)\otimes1=\bigwedge^n(u\otimes1)$, donc on peut supposer
$n=1$. Mais alors l'application $u\mapsto\det u$ est un homomorphisme du
$A$-module $\operatorname{Hom}_A(E,E)$ dans $K$~; comme $E$ est de type fini,
$\operatorname{Hom}_A(E,E)$ est isomorphe à un sous-module du $A$-module de
type fini $E^n$ (si $E$ admet un système de $n$ générateurs)~; donc les
éléments $\det u$ appartiennent à un sous-$A$-module de $K$ de type fini, et
sont par suite entiers sur $A$.

